\section{DDIM Inversion：确定性反演与图像编辑}

\subsection{DDIM 的确定性采样回顾}

DDIM（Denoising Diffusion Implicit Models）的核心特性是将反向过程变为\textbf{确定性}的。在 DDPM 中，反向过程的每一步都注入随机噪声；而在 DDIM 中，设方差 $\sigma_t = 0$，则更新公式变为：

$$
x_{t-1} = \sqrt{\bar\alpha_{t-1}}\,\hat{x}_0(x_t) + \sqrt{1-\bar\alpha_{t-1}}\,\epsilon_\theta(x_t, t)
$$

其中 $\hat{x}_0(x_t) = \frac{x_t - \sqrt{1-\bar\alpha_t}\,\epsilon_\theta(x_t, t)}{\sqrt{\bar\alpha_t}}$ 是对干净图像的估计。

\begin{importantbox}{DDIM 确定性性质}
在 DDIM 中，从 $x_T$ 到 $x_0$ 的映射是确定性的：给定相同的初始噪声 $x_T$ 和模型参数，总是生成完全相同的 $x_0$。这意味着存在一个隐式的\textbf{双射映射} $x_T \leftrightarrow x_0$。
\end{importantbox}

\subsection{反演：从 $x_0$ 找回 $x_T$}

DDIM 的确定性特性使得一个重要操作成为可能——\textbf{反演}（inversion）。给定一张真实图像 $x_0$，我们希望找到对应的噪声向量 $x_T$，使得 DDIM 的确定性反向过程能精确重建 $x_0$。

\begin{warningbox}{反演并非简单的前向过程}
直觉上，前向过程 $q(x_T|x_0) = \mathcal{N}(\sqrt{\bar\alpha_T}\,x_0,\, (1-\bar\alpha_T)\,I)$ 可以直接计算 $x_T$。但这是\textbf{随机}的前向过程，得到的 $x_T$ 并不保证通过 DDIM 确定性反向过程回到 $x_0$。DDIM Inversion 需要沿着\textbf{确定性}路径反向计算。
\end{warningbox}

反演的近似方法：将 DDIM 更新公式中的 $t-1$ 替换为 $t+1$，利用噪声函数的平滑性假设（$\epsilon_\theta(x_t, t) \approx \epsilon_\theta(x_{t+1}, t+1)$），逐步从 $x_0$ 计算到 $x_T$：

$$
x_{t+1} = \sqrt{\bar\alpha_{t+1}}\,\hat{x}_0(x_t) + \sqrt{1-\bar\alpha_{t+1}}\,\epsilon_\theta(x_t, t)
$$

\begin{knowledgebox}{近似精度与时间步数}
DDIM Inversion 的精度取决于时间步间隔的大小。时间步数越多（间隔越小），$\epsilon_\theta(x_t, t) \approx \epsilon_\theta(x_{t+1}, t+1)$ 的近似越准确，反演重建越精确。实践中通常需要 50--200 步才能获得高质量的反演结果。
\end{knowledgebox}

\subsection{DDIM Inversion + CFG = 图像编辑}

DDIM Inversion 与 CFG 结合可以实现强大的图像编辑功能。流程如下：

\begin{enumerate}
    \item \textbf{反演阶段}：给定输入图像 $x_0$ 和原始文本描述 $y_{\text{src}}$（如 ``a photograph of a puppy''），通过 DDIM Inversion 得到 $x_T$
    \item \textbf{生成阶段}：从相同的 $x_T$ 出发，使用新的文本描述 $y_{\text{tgt}}$（如 ``a photograph of a gray cat''）进行 DDIM 确定性采样
    \item \textbf{结果}：输出图像在整体构图和风格上与原图相似，但内容符合新的文本描述
\end{enumerate}

这种方法可以实现：
\begin{itemize}
    \item 物体替换：将水果替换为蛋糕
    \item 属性修改：给人物戴上太阳镜
    \item 风格变换：将照片变为特定绘画风格
\end{itemize}

\subsection{本章小结}

DDIM Inversion 利用 DDIM 采样的确定性特性，实现了从真实图像到噪声空间的近似映射。与 CFG 结合后，只需更换文本条件即可实现语义级别的图像编辑，而无需任何额外训练。这展示了扩散模型不仅能生成图像，还能作为强大的图像编辑工具。

\newpage
